
\section{Introduction}
\label{sec:introduction}

A core challenge in deploying language models (LMs) is \textit{verification}: determining the quality or correctness of a model's response. 
This problem arises across various components of the LM  pipeline, including dataset curation, model alignment, and inference-time decision-making.  
Verification relies on \textit{verifiers}—functions that score responses. 
When combined with repeated sampling---generating multiple candidate responses from a LM---a perfect verifier can be used to select a correct candidate response, significantly enhancing model capability on tasks such as math, code, and reasoning \citep{snell2024scalingllmtesttimecompute, brown2024largelanguagemonkeysscaling, puri2025probabilisticinferenceapproachinferencetime}. 
For example, Llama 3.1 8B Instruct can match Llama 3.1 70B Instruct and even GPT-4o performances on MATH500 \citep{hendrycks2021measuring} and MiniF2F \citep{zheng2022minif2fcrosssystembenchmarkformal} when paired with perfect verifiers for these mathematics tasks.  
However, without a perfect verifier, a \textit{generation-verification gap} emerges \citep{song2025mindgapexaminingselfimprovement}: a LM can generate a correct response, but we fail to identify it. 

The generation-verification gap is prevalent across many tasks across mathematics, coding, scientific reasoning, instruction-following, and more.
For some of these settings, we have access to \textit{oracle verifiers} that can perfectly identify correct responses. A prominent example is Lean, a formal theorem prover that can be used for problems such MiniF2F \citep{zheng2022minif2fcrosssystembenchmarkformal}. However, this is often a limited setup, as not all mathematical proofs can be processed by Lean. 
Alternatively, humans could judge LM responses but manual evaluation is often expensive, noisy, and difficult to scale \citep{hosking2024humanfeedbackgoldstandard, clark2021thatshumangoldevaluating, karpinska2021perilsusingmechanicalturk}.
In contrast, LMs prompted as judges \citep{chiang2024chatbot} and reward models \citep{lambert2024rewardbenchevaluatingrewardmodels, singhi2025solveverifycomputeoptimalproblem, liu20251bllmsurpass405b} can be applied off-the-shelf to tasks like mathematics, coding, scientific reasoning, instruction-following \citep{hendrycks2021measuring, rein2024gpqa, jain2024livecodebenchholisticcontaminationfree, alpaca_eval}. 
However, these \textit{weak verifiers} produce noisy, inconsistent scores, often exhibit poor calibration, and suffer from high false positive rates \citep{stroebl2024inferencescalingflawslimits}. 
We ask: \textit{to what extent can we leverage weak verifiers to improve accuracy in the repeated sampling regime?}

We explore \textit{scaling verification}, specifically how to combine \textit{multiple weak verifiers} to improve response selection for repeated sampling.
As new pre-trained models become available, the pool of weak verifiers continues to expand and offer diverse, complementary sources of signal that could improve response selection if they can be aggregated effectively.
Recent work has explored scaling verification through techniques such as self-verification or averaging LM judge scores \citep{lifshitz2025multiagentverificationscalingtesttime, zhao2025sample, chen2025sets} although other work has found limitations to scaling test-time compute when utilizing weak verifiers for response selection \citep{stroebl2024inferencescalingflawslimits}.  
We observe three key challenges towards ensembling weak verifiers:

\input{tables_and_figures/main_figure}

\begin{enumerate} [itemsep=0.00pt,topsep=0pt,leftmargin=12pt]
    \item \textbf{Naively aggregating weak verifiers is insufficient for reliable verification.} Weak verifiers such as LM-based judges or reward models produce noisy, biased, and poorly calibrated scores, leading to inconsistent performance. \citep{stroebl2024inferencescalingflawslimits, lambert2024rewardbenchevaluatingrewardmodels, chiang2024chatbot}. 
    While using a naive unweighted average of verifier scores is straightforward, it implicitly assumes uniform verifier quality, causing low-quality verifiers to dominate and degrade the overall accuracy~\citep{verga2024replacingjudgesjuriesevaluating, xu2024perfectblendredefiningrlhf, eisenstein2023helping}. Moreover, while previous work has hypothesized that more sophisticated weighted ensembles should perform better, this claim has not been studied~\citep{lifshitz2025multiagentverificationscalingtesttime}. 
   
    \item \textbf{Effective ensembling with limited labeled data is challenging.} More sophisticated ensembling techniques typically learn verifier weights from labeled data, but such data is expensive and difficult to obtain. 
    \textit{Weak Supervision} (WS), a family of statistical techniques developed for data labeling, offers a potential solution through algorithms that aggregate multiple weak signals---such as crowd-worker annotations and expert-defined heuristics---while only requiring a small amount of labeled data \citep{ratner2016data, ratner2019training, fu2020fast}.  
    In traditional WS, practitioners can design and shape each weak signal to ensure sufficient quality (i.e., iteratively tweaking program-based heuristics), and guarantees of WS hinge on a baseline level of quality.
    Our weak signals, however, are fixed pre-trained language model verifiers, which have wildly varying accuracy---especially when applied to out-of-distribution tasks---and can emit incompatible outputs (logits, binary scores, Likert scores)~\citep{lambert2024rewardbenchevaluatingrewardmodels} that we cannot easily tweak. Due to these conditions, WS algorithms may not perform well when directly applied to verification.
    \item \textbf{Verification is expensive to deploy at inference.}
    Verification can dominate inference-time costs \citep{singhi2025solveverifycomputeoptimalproblem, liu20251bllmsurpass405b}, since each verifier must process both the problem and its candidate response(s) \cite{lightman2023letsverifystepstep}, often evaluating intermediate steps \cite{lightman2023letsverifystepstep} and multiple solution paths \cite{snell2024scalingllmtesttimecompute}. 
    In fact, achieving gains over unverified generation (i.e. majority voting) can require $10\times$ to $128\times$ the inference compute per query \citep{singhi2025whentosolve, lifshitz2025multiagentverificationscalingtesttime, zhao2025sample, chen2025sets}. 
    
\end{enumerate}


In this work, we introduce \weaver{}, a framework for aggregating weak verifiers without supervised finetuning on ground truth labels (Figure \ref{fig:main_figure}). 
First, we demonstrate that if we have access to a large corpus of labeled training data (e.g., 50,000 query-response pairs), we can learn weighted ensembles that can outperform naive averaging by up to 11.2\% points. This is because weighted ensembles take advantage of wide variability in verifier accuracy. However, in many real-world scenarios, we do not have access to such quantities of labeled data.
Second, to reduce the dependency on labeled data, we adapt Weak Supervision to the verification setting by addressing challenges around inconsistent outputs and low-accuracy verifiers. 
\weaver{} filters out uninformative verifiers, normalizes verifier scores, and builds a latent variable model over these scores and the unknown true labels to estimate the verifier accuracies to be used as weights for the ensemble~\citep{ratner2016data, hall2003generalized}.

Empirically, given a repeated sampling budget and a set of verifiers, \weaver{} improves over repeated sampling with unweighted averaging of verifier scores by 17.1\% and with majority voting by 13.5\% (Table \ref{tab:verifier_ablations}; Figure \ref{fig:fp_and_pass@1}).
Compared to an LM's $Pass@1$, \weaver{} allows us to improve performance by 17.9\% for 8B models and 14.5\% for 70B models across reasoning and mathematics tasks (Tables \ref{tab:verifier_ablations} and \ref{tab:8B_verifier_ablations}).
\textit{This mirrors the performance jump from GPT-4o to o3-mini} (73.9\% vs. 88.2\%)---but only via increased sampling at test time rather than parameter tuning or post-training procedures. 
We also study how \weaver{} scales along different axes of test-time compute: generation, verifiers, model size, and inference budget (Section~\ref{sec:scaling_verification_with_weaver}). 
We find that even as we increase the number of generations, many standard verification baselines (e.g. majority voting) quickly plateau (Figure \ref{fig:fp_and_pass@1}). 
Naive ensembling saturates more slowly, but its gains are limited by sensitivity to the model choice and the number of verifiers. 

Finally, to mitigate the compute costs of calling multiple weak verifiers for each response, we extend \weaver{} by training a 400M-parameter cross-encoder verifier using \weaver{}'s selected responses.
We demonstrate that using a distilled \weaver{} cross-encoder as a verifier \textit{retains 98.7\% of the accuracy gains} from the learned verifier ensemble  while reducing compute costs by three orders of magnitude -- \textit{saving 99.97\% inference FLOPS} while still capturing an effective verification strategy (Section \ref{sec:weaver_distillation}). 
Overall, our findings highlight that more reliable, scalable verification is possible even in the absence of ground-truth labels---paving the way for improved data filtering, model alignment, and inference-time decision-making. 
